%%%%%%%%%%%%%%%%%%%%%%%%%%%%%%%%%%%%%%%%%%%%%%%%%%%%%%%%%%%%%%%%%%%%%%%%%%%%%%%
%% Apêndice B -- Roteiro para elaboração dos relatórios
%%%%%%%%%%%%%%%%%%%%%%%%%%%%%%%%%%%%%%%%%%%%%%%%%%%%%%%%%%%%%%%%%%%%%%%%%%%%%%%
\apendice{Roteiro para Elaboração dos Relatórios}

\subtituloguia{1) Instruções gerais}

O relatório deverá ser feito à mão. Erros no uso correto do Português são penalizados
com redução de nota. A parte inicial do relatório deverá conter:

\begin{enumerate}[label=\alph*),leftmargin=1.2cm]
  \item nome e número de matrícula dos alunos que realizaram a prática;
  \item o número e o título da prática;
  \item a data em que se realizou a prática.
\end{enumerate}

\subtituloguia{2) Introdução teórica}

Esta seção deverá conter uma introdução sobre a prática realizada, com enfoque no
objetivo e no conteúdo que foi avaliado experimentalmente. Expresse os principais
circuitos relacionados à prática, bem como os dados obtidos por simulação e/ou
experimentalmente e os calculados. Liste todos os materiais utilizados na prática.

\subtituloguia{3) Memória de cálculo}

Nesta seção deverão ser apresentadas as questões, tabelas e/ou gráficos relativos à
parte Prática de cada um dos experimentos. É permitido responder no próprio roteiro da
aula prática e anexá-lo ao relatório.

\subtituloguia{4) Resultado da pesquisa}

Nesta seção deverá ser apresentado o assunto pesquisado, caso exista esse tópico no
roteiro da prática realizada. Não se esqueça de colocar a fonte da pesquisa (livro,
internet, artigos).

\subtituloguia{5) Conclusão}

Nesta seção deve-se elaborar e redigir as principais conclusões observadas pelo grupo.
A conclusão deve versar sobre os objetivos. É importante ressaltar que a conclusão é a
parte mais importante do relatório.
